\chapter{Fluid from the Ear}

\section*{Definition and Overview}
Fluid from the ear, also called ear discharge or otorrhoea, is any liquid that comes out of the ear canal. It can be watery, bloody, thick, or pus-like, and it has a range of causes, from harmless earwax and moisture to infections of the outer or middle ear, and, rarely, serious conditions such as a perforated eardrum, a foreign body, or injury. In many cases the discharge is accompanied by pain, itching, or hearing loss, and the pattern of the discharge gives clues to its cause. Most causes of ear discharge are treatable, but some, including a perforated eardrum or discharge after head injury, need urgent medical attention. This chapter explains the causes of fluid from the ear and what to do about it.

\section*{Epidemiology}
Ear discharge is a common presentation in general practice and emergency departments, particularly in children. Outer ear infections (otitis externa, or swimmer's ear) are common in swimmers and in hot, humid climates, and middle ear infections with a burst eardrum cause discharge mainly in young children. A perforated eardrum from infection, injury, or pressure change discharges pus, blood, or clear fluid. While most causes are minor and resolve with treatment, chronic discharge can indicate persistent infection or a cholesteatoma, a serious condition, and discharge after head injury can signal a skull fracture and requires emergency care.

\section*{Aetiology and Pathophysiology}
\begin{itemize}
    \item \textbf{Earwax:} normal cerumen can melt, appear watery, or be mistaken for discharge
    \item \textbf{Outer ear infection (otitis externa):} bacterial or fungal infection of the ear canal, often from swimming, moisture, scratching, or foreign objects; the canal swells, itches, and discharges pus or fluid
    \item \textbf{Middle ear infection with perforation:} acute otitis media builds pressure behind the eardrum; if it bursts, pus drains out, often relieving the pain
    \item \textbf{Perforated eardrum:} from infection, injury (including cotton buds), or sudden pressure changes, with blood-stained or clear discharge
    \item \textbf{Foreign bodies:} objects in the ear canal, common in children, cause discharge as the canal becomes inflamed
    \item \textbf{Cholesteatoma:} a destructive growth of skin in the middle ear that causes foul-smelling, persistent discharge and progressive hearing loss
    \item \textbf{Chronic suppurative otitis media:} persistent middle ear infection with a permanent perforation and ongoing discharge
    \item \textbf{Head injury:} clear fluid (cerebrospinal fluid) leaking from the ear after head trauma is an emergency
    \item \textbf{Rare causes:} tumours of the ear canal or middle ear
\end{itemize}

\section*{Clinical Features}
\begin{itemize}
    \item The appearance of the discharge: watery, sticky pus, blood-stained, or clear and colourless
    \item Pain, which may be severe in outer or middle ear infection and often eases when the eardrum perforates and drains
    \item Itching in the ear canal, typical of otitis externa
    \item Hearing loss, fullness, or ringing in the affected ear
    \item Fever, especially with middle ear infection
    \item Recent swimming, ear cleaning with cotton buds, a foreign body, or head injury in the history
    \item Foul smell, which suggests cholesteatoma or chronic infection
    \item Clear fluid leaking from the ear after head injury, which is a medical emergency
\end{itemize}

\section*{Differential Diagnosis}
The causes are distinguished by examination and history.
\begin{itemize}
    \item Otitis externa versus acute otitis media with perforation
    \item Perforated eardrum from injury or infection
    \item Foreign body, common in young children
    \item Cholesteatoma, with persistent foul discharge and hearing loss
    \item Chronic suppurative otitis media
    \item Wax and water in the ear, which are not infections
    \item Cerebrospinal fluid leak after head injury, which is an emergency
\end{itemize}

\section*{When to Seek Medical Care}
Most ear discharge should be assessed by a doctor, especially if it is painful, persistent, associated with hearing loss or fever, or occurring in a child. Discharge after head injury needs urgent assessment.

\textbf{Call 000 (or local emergency services) for:}
\begin{itemize}
    \item Clear or bloody fluid leaking from the ear after a head injury
    \item Sudden severe ear or head pain
    \item Fever, confusion, or a stiff neck with ear discharge, which may indicate spread of infection
    \item Facial weakness or severe dizziness with ear symptoms
    \item Collapse or altered consciousness
\end{itemize}

\section*{Diagnosis and Investigations}
\begin{itemize}
    \item \textbf{History:} the appearance and duration of discharge, pain, hearing, recent swimming or injury, and foreign body risk
    \item \textbf{Examination:} the ear canal and eardrum are examined with an otoscope; gentle cleaning may be needed to see the eardrum
    \item \textbf{Swab of the discharge:} to identify the organism and guide antibiotic choice
    \item \textbf{Hearing tests:} for persistent or recurrent discharge
    \item \textbf{Microsuction:} safe cleaning of the ear by a specialist when the canal is blocked
    \item \textbf{CT or MRI:} for suspected cholesteatoma or complications
    \item \textbf{Imaging after head injury:} for suspected skull fracture
\end{itemize}

\section*{Management}
\begin{itemize}
    \item \textbf{Outer ear infection:} antibiotic or antifungal ear drops, keeping the ear dry, and avoiding further irritation; wicks may be used if the canal is very swollen
    \item \textbf{Middle ear infection with perforation:} oral antibiotics when bacterial, with analgesia; the eardrum usually heals itself
    \item \textbf{Perforated eardrum:} keep the ear dry, avoid water and cotton buds, treat infection, and review healing; some perforations need surgical repair
    \item \textbf{Foreign body:} removal by a doctor; never attempt to dig it out at home
    \item \textbf{Cholesteatoma:} surgical removal, because it destroys bone and hearing if untreated
    \item \textbf{Chronic discharge:} aural toilet, topical treatment, and specialist referral
    \item \textbf{After head injury:} emergency assessment, including imaging, and the person should not have anything inserted in the ear
    \item \textbf{Pain relief:} paracetamol or ibuprofen
    \item \textbf{Ear protection:} keep water out of the ear while the eardrum is perforated
\end{itemize}

\section*{Complications}
\begin{itemize}
    \item Hearing loss, temporary or permanent
    \item Spread of infection to the mastoid bone (mastoiditis) or the brain (meningitis, brain abscess), in severe cases
    \item Facial nerve weakness from cholesteatoma or severe infection
    \item Chronic discharge and permanent perforation
    \item Damage to the ear canal or eardrum from inserting objects
    \item Injury to the eardrum from attempted home removal of foreign bodies
\end{itemize}

\section*{Prognosis and Outcomes}
Most ear discharge resolves completely with appropriate treatment. Acute otitis media with perforation usually heals within weeks with normal hearing, and otitis externa clears with drops and dry-ear measures. Chronic conditions such as cholesteatoma have an excellent outcome when detected and treated surgically, but progressive damage if ignored. The key to good outcomes is prompt assessment, accurate diagnosis, and avoiding the instinct to insert anything into the ear, which aggravates most causes.

\section*{Prevention and Self-Care}
\begin{itemize}
    \item Never insert cotton buds, fingers, or other objects into the ear
    \item Dry the ears after swimming with a towel, and use a hair dryer on a cool setting at a distance if prone to swimmer's ear
    \item Wear earplugs when swimming if prone to outer ear infections
    \item Treat ear infections promptly to prevent perforation
    \item Keep the ear dry when the eardrum is perforated
    \item Remove foreign bodies only with medical help
    \item See a doctor for discharge that persists, smells, or is associated with hearing loss
    \item Seek urgent care for discharge after head injury
\end{itemize}

\section*{Sources and Further Reading}
The following reputable sources provide further information for healthcare students and patients seeking reliable, current advice about fluid from the ear.
\begin{itemize}
    \item Healthdirect Australia. \textit{Ear discharge.} https://www.healthdirect.gov.au/ear-discharge
    \item Better Health Channel. \textit{Otitis externa and ear discharge.} https://www.betterhealth.vic.gov.au/
    \item The Royal Children's Hospital Melbourne. \textit{Ear discharge fact sheet.} https://www.rch.org.au/
    \item NHS. \textit{Ear discharge.} https://www.nhs.uk/conditions/ear-discharge/
    \item Mayo Clinic. \textit{Ear discharge.} https://www.mayoclinic.org/
    \item MedlinePlus. \textit{Ear discharge.} https://medlineplus.gov/
\end{itemize}
